% answers_short.tex —— 2015 年·广东卷·文科 · 简短答案速查
% 由 tools/make_short_answers.py 自 answers.tex 的 \gksolhead/\gkanswer 提取，
% 并与 2015/2015_数学真题_广东文_答案要点.md 逐题比对后生成，勿手改（改源 MD 后重跑）。

\begin{gknotice}
\textbf{使用说明}：本节只给最终答案便于快速核对，完整推导见解析版。答案为编写者独立推导，\textbf{非官方标准答案}。
\end{gknotice}

\gkpart{选择题}
\begin{center}\begin{tabular}{ccccccccccc}
\toprule
题号 & 1 & 2 & 3 & 4 & 5 & 6 & 7 & 8 & 9 & 10 \\
\midrule
答案 & B & A & D & B & C & D & B & B & A & A \\
\bottomrule
\end{tabular}\end{center}

\gkpart{填空题}
\begin{itemize}[leftmargin=2.6em]
  \item[\textbf{11.}] \((-4,1)\)
  \item[\textbf{12.}] \(11\)
  \item[\textbf{13.}] \(1\)
  \item[\textbf{14.}] \((2,-4)\)
  \item[\textbf{15.}] \(3\)
\end{itemize}

\gkpart{解答题}
\begin{itemize}[leftmargin=2.6em]
  \item[\textbf{16.}] (1) \(-3\)；(2) \(1\)
  \item[\textbf{17.}] (1) \(x=0.0075\)；(2) 众数 \(230\)（度），中位数 \(224\)（度）；(3) \(5\) 户
  \item[\textbf{18.}] (1) 成立（长方形 \(ABCD\) 中 \(BC\parallel AD\)，\(AD\subset\) 平面 \(PDA\)，\(BC\not\subset\) 平面 \(PDA\)，故 \(BC\parallel\) 平面 \(PDA\)）；(2) 成立（平面 \(PDC\perp\) 平面 \(ABCD\)，交线为 \(DC\)；\(BC\perp DC\)，故 \(BC\perp\) 平面 \(PDC\)，从而 \(BC\perp PD\)）；(3) \(\dfrac{3\sqrt{7}}{2}\)
  \item[\textbf{19.}] (1) \(a_4=\dfrac{7}{8}\)；(2) 成立（\(\left\{a_{n+1}-\dfrac{1}{2}a_n\right\}\) 是首项为 \(1\)、公比为 \(\dfrac{1}{2}\) 的等比数列）；(3) \(a_n=\dfrac{2n-1}{2^{n-1}}\)
  \item[\textbf{20.}] (1) \((3,0)\)；(2) \(\left(x-\dfrac{3}{2}\right)^2+y^2=\dfrac{9}{4}\)（\(\dfrac{5}{3}<x\leqslant 3\)）；(3) 存在，\(k\in\left[-\dfrac{2\sqrt{5}}{7},\dfrac{2\sqrt{5}}{7}\right]\cup\left\{-\dfrac{3}{4},\dfrac{3}{4}\right\}\)
  \item[\textbf{21.}] (1) \(a\leqslant\dfrac{1}{2}\)；(2) \(f(x)\) 在 \((-\infty,a)\) 上单调递减，在 \((a,+\infty)\) 上单调递增；(3) 当 \(a=2\) 时，有 \(1\) 个零点；当 \(a>2\) 时，有 \(2\) 个零点
\end{itemize}
